\section{Local coefficient identities}\label{app:local-calculations}
This appendix proves Lemma~\ref{lem:local-identities} in the diagram convention of Section~\ref{sec:flux}. The finite tables also specify every nonzero local coefficient used by the interpolation argument.

\begin{proof}
Equation \eqref{eq:01-factorization} follows by direct substitution.
For completeness we record a coefficient check of the splitting
identity. Subscripts $1,2$ mean arguments $x+s\lambda,x$. Let $q$ be
the cup coefficient and $k$ the strand coefficient of the splitter:
$q=k=\rho$ for $s=2$, and $q=1,k=0$ for $s=3$. Unsubscripted weights
refer to $\widehat R_s$; in the deletion case use $B=1$ and set its
other nonempty two-slot coefficients to zero. The inputs are the
passing auxiliary strand and the slot to be split; the latter is
absent in the deletion case. Outputs $1,2$ are the split slots and
output $3$ is the auxiliary slot. Input bits record occupancy. For one
input strand, a bare output index specifies its destination and a
parenthesized pair specifies an additional cup. The complete coefficient
check is
\[
\begin{array}{c|c|c|c}
\text{input}&\text{connections at output}&\text{left side}&(S_s\otimes I)\widehat R\\\hline
00&\emptyset&1&1\\
&12\ (\text{cup})& U_1 A_2+q B_1 B_2 &q\\
&13&U_1 B_2+q B_1 A_2 & k U\\
&23&U_2+q A_1 E_2&k U\\\hline
10&1&A_1+q F_1 U_2 & k A\\
&2&B_1 A_2+q U_1 B_2 &k A\\
&3&B_1 B_2+q U_1 A_2 & B\\
&3,(12)&q F_1 E_2 &q B\\
&1,(23)&A_1 U_2+q(E_1 E_2+F_1 F_2)&0
\end{array}
\]
\begin{samepage}
For $s=2$ there are also the following inputs. In the final $11$ block,
an empty output or a parenthesized output pair connects the two input
strands to one another; a bare pair of output indices denotes two
through-strands.
\[
\begin{array}{c|c|c|c}
01&1&k(B_1+U_1 U_2)& k B\\
&2&k(A_1 A_2+B_2)& k B\\
&3&k(A_1 B_2+A_2)& A\\
&1,(23)& k(B_1 U_2+U_1 F_2)&0\\
&3,(12)& k U_1 E_2 & q A\\\hline
11&\emptyset&k(U_1+B_1 U_2)& U\\
&12& k(E_1 A_2+A_1 B_2)&0\\
&(12)& k F_1 A_2 &q U\\
&13& k(E_1 B_2+A_1 A_2)& k E\\
&(13)& k F_1 B_2&k F\\
&23& k B_1 E_2&k E\\
&(23)&k(U_1 U_2+B_1 F_2)& k F
\end{array}
\]
\end{samepage}
Here are the substitutions verifying these equalities. Put
$s_j=\sin(x+j\lambda)$. The weights at $x+j\lambda$ have common
denominator $s_{j+2}s_{j+3}$ and numerators
\[
\begin{array}{c|ccccc}
\text{weight}&A&B&U&E&F\\\hline
\text{numerator}&d s_{j+5}&s_j s_{j+5}&d s_j&
 s_{j+5}s_{j+6}&-s_{j+7}s_j
\end{array}
\]
Substitution uses the following product-to-sum identities:
\[
\begin{gathered}
s_{j+8}=-s_j,\quad s_j s_{j+1}+s_{j+4}s_{j+5}=C,\quad
s_2s_6=d^2-s_0^2,\quad s_1s_5=d^2-s_7^2,\\
\rho(s_1-s_7)=s_0,\quad \rho(s_6-s_0)=s_7,\quad
\rho(s_1s_2-s_6s_7)=s_0s_4,\quad \rho(s_5s_6-s_0s_1)=s_7s_3,\\
\rho(s_7s_0+s_4s_5)=s_6s_2,\quad \rho(s_2s_3+s_7s_0)=s_1s_5,\quad
s_1s_6-s_7s_0=\rho d^2,\\
s_1s_6s_5=s_7s_2s_3+d^2s_0,\qquad s_1s_6s_2=s_0s_4s_5+d^2s_7
\end{gathered}
\]
The last two trigonometric identities also follow from the preceding
formulas for $s_1s_5$ and $s_6s_2$. Thus substitution verifies every
entry in the tables.

For \eqref{eq:01-unitarity}, the sector with one strand reduces to
multiplication of two matrices with diagonal entries $A$ and
off-diagonal entries $B$. In the even sector the required identities
are
\[
 E(x)E(-x)=1,\qquad U(x)+E(x)U(-x)=0,
\]
\[
 E(x)F(-x)+F(x)E(-x)+U(x)U(-x)=0.
\]
A cap composed with a cup creates a loop and therefore vanishes.
These identities again follow from \eqref{eq:01-rhombic-weights}.

To prove \eqref{eq:01-braid}, first set $y=s\lambda$, $s=2,3$.
Use \eqref{eq:01-factorization} and \eqref{eq:01-splitting}, together
with their transposes and the reversal of slot order. The effective
box is symmetric, and both sides become
$(I\otimes S_s)\widehat R_s(x)(S_s^{\mathsf t}\otimes I)$.
Unitarity transports these identities to $x+y=s\lambda$, by applying
the identity with arguments $(-x,x+y)$. The case $y=0$ follows from
$R(0)=I$.

Multiply the difference of the two sides by their three scalar
denominators, and call the resulting diagram sum $\mathcal F(x,y)$.
Each coefficient is a Laurent polynomial in $e^{iy}$ with exponents
between $-4$ and $4$. To see its parity, let $D_j$ multiply a diagram
by $-1$ when slot $j$ is occupied. Inspection of the nine boxes gives
\[
 R_i(v+\pi)=D_{i+1}R_i(v)D_i.
\]
On either side of the braid relation the two internal $D_2$ signs
cancel. The remaining signs are at input slot $1$ and output slot $3$;
the scalar denominators are $\pi$-periodic. Hence
\[
 \mathcal F(x,y+\pi)=D_3\mathcal F(x,y)D_1.
\]
Every boundary coefficient has fixed parity. After multiplication by
a Laurent monomial it is a polynomial of degree at most four in
$e^{2iy}$. We have proved its vanishing at
$y=0,2\lambda,3\lambda,2\lambda-x,3\lambda-x$; these five values
are distinct modulo $\pi$ for generic $x$. Interpolation proves the
identity, and rational continuation removes the genericity restriction.
\end{proof}
